(1) 有效两模式体系以 \(H=\hbar\begin{pmatrix}i\gamma&J\\J&-i\gamma\end{pmatrix}\) 描述平衡增益损耗，\(J>0,\gamma\ge0\) 为频率量纲，标准内积为普通复向量Euclid内积。求本征能量，分别讨论 \(0\le\gamma<J\)、\(\gamma=J\)、\(\gamma>J\) 的谱及 \(\gamma=J\) 的本征向量并合，给出PT对称性的含义。
(2) 由 \(H^2\) 构造任意实时间的精确传播子 \(U(t)=e^{-iHt/\hbar}\)；特别写出例外点的多项式传播子。对初态 \(|L\rangle=(1,0)^T\)，在未破缺区求两个振幅及普通范数 \(N(t)\)，验证它通常不守恒。
(3) 计算普通范数的导数，证明它等于 \(2\gamma(|a_L|^2-|a_R|^2)\)，并解释非厄米有效哈密顿量不应被解释成一个孤立系统的封闭Hermitian概率演化。
(4) 在 \(0\le\gamma<J\) 区间构造Hermitian正定的 \(2\times2\) 度量 \(\eta\)，满足 \(H^\dagger\eta=\eta H\)。求度量特征值，证明 \(\langle\psi|\eta|\psi\rangle\) 守恒，解释为何正定性在例外点处丧失及 \(\gamma\to0\) 的幺正极限。
